\section{Worum es geht}

\textbf{[SETZUNG]} In der FFGFT folgt $\alpha$ aus der Geometrie (A130).
Dieses Dokument behandelt die komplementäre Frage: wie setzt man $\alpha=1$
\emph{als Rechenwerkzeug} in einer Weise, die vollständig reversibel ist
und die Kopplungsstruktur sichtbar macht? Die Antwort ist die
Heaviside-Lorentz-Setzung $e^2=4\pi\alpha$ kombiniert mit $\alpha=1$ --
keine Einheitenwahl, sondern eine \emph{Stipulation} mit erhaltener Spur.

\section{Setzung: Heaviside-Lorentz mit $\alpha = 1$}

\textbf{[SETZUNG]} In natürlichen Einheiten ($\hbar=c=1$) und im
Heaviside-Lorentz-System gilt als \emph{Definitionsgleichung}
\[
  e^2 = 4\pi\alpha.
\]
Die Feinstrukturkonstante wird auf $\alpha=1$ gesetzt. Das ist keine
\emph{Norm} (Einheitenwahl) -- eine dimensionslose Zahl ist gegen
Umskalierung der Basiseinheiten invariant --, sondern eine
\emph{Stipulation}: die Kopplung wird als Rechen-Einheit adoptiert.

\textbf{[B]} Die Ladung trägt das $4\pi$; auf $1$ gesetzt wird $\alpha$,
nicht $e$. Aus $e^2=4\pi\alpha$ folgt mit $\alpha=1$ der Wert
$e=\sqrt{4\pi}\approx3{,}545$ -- nicht $e=1$. Die Relation $e^2=4\pi\alpha$
bleibt als Definitionsgleichung stehen und kann nicht durch die Setzung
aufgelöst werden.

\section{Rückrechnung als Substitution}

\textbf{[B]} Weil $e^2=4\pi\alpha$ als Relation stehen bleibt, ist die
Kopplungspotenz automatisch mitgeführt. Die Rückrechnung ist reine Substitution
\[
  \alpha \;\longrightarrow\; \frac{1}{137{,}036},
\]
angewandt auf die $\alpha$-Ordnung des Prozesses:
\begin{itemize}[leftmargin=2em,itemsep=2pt,topsep=2pt]
  \item pro Vertex ein Faktor $e=\sqrt{4\pi\alpha}$; gegenüber dem
    $\alpha{=}1$-Wert $\sqrt{4\pi}$ also ein $\sqrt\alpha$ zurück,
  \item Amplitude mit $V$ Vertizes $\propto(4\pi\alpha)^{V/2}
    \;\Rightarrow\;\times\alpha^{V/2}$,
  \item Rate / Wirkungsquerschnitt $\propto(4\pi\alpha)^V
    \;\Rightarrow\;\times\alpha^V$,
  \item kompakt: Prozess der Rate-Ordnung $\alpha^n
    \;\Rightarrow\;\times\alpha^n$.
\end{itemize}

\textbf{[B]} Das $4\pi$ wird dabei nicht doppelt gezählt: es steckt bereits
in $e^2$ und in der Phasenraum-Normierung. Zurückgesetzt wird nur $\alpha^n$,
nicht $(4\pi\alpha)^n$.

\section{Umrechnungsbeispiele}

\textbf{[K]} Alle Größen in $\hbar=c=1$, Heaviside-Lorentz.
Rückfaktor $=$ die $\alpha$-Ordnung des Ausdrucks.

\begin{center}
\renewcommand{\arraystretch}{1.35}
{\small
\begin{tabular}{lllc}
\toprule
Größe & physikalische Form & mit $\alpha{=}1$ & Rückfaktor \\
\midrule
Coulomb-Potential      & $V = \alpha/r$                         & $1/r$           & $\times\alpha$      \\
Klass.\ Elektronradius & $r_e = \alpha/m_e$                     & $1/m_e$         & $\times\alpha$      \\
Compton-Wellenlänge    & $\lambda_C = 1/m_e$                    & $1/m_e$         & $\times 1$          \\
Bohr-Radius            & $a_0 = 1/(\alpha m_e)$                 & $1/m_e$         & $\times\alpha^{-1}$ \\
Rydberg / Bindung      & $E = \tfrac12\alpha^2 m_e$             & $\tfrac12 m_e$  & $\times\alpha^2$    \\
Thomson-Querschnitt    & $\sigma_T = \tfrac{8\pi}{3}\alpha^2/m_e^2$ & $\tfrac{8\pi}{3}/m_e^2$ & $\times\alpha^2$ \\
\bottomrule
\end{tabular}
}
\renewcommand{\arraystretch}{1.0}
\end{center}

\section{Die 137-Leiter: was $\alpha=1$ sichtbar macht}

\textbf{[B]} Die drei atomaren Längenskalen stehen im Verhältnis
\[
  r_e : \lambda_C : a_0 \;=\; \alpha : 1 : \alpha^{-1}.
\]
Bei $\alpha=1$ fallen sie alle auf $1/m_e$ zusammen: sie sind keine drei
unabhängigen Skalen, sondern \emph{eine} Skala mit drei $\alpha$-Potenzen.
$a_0$ ist genau $137\times$ größer als $\lambda_C$, und $\lambda_C$ genau
$137\times$ größer als $r_e$.

\textbf{[B]} Mit $\alpha=1/137$ überall ausgeschrieben bleibt diese gemeinsame
Herkunft verdeckt -- da steht dreimal eine andere krumme Zahl und die
$\alpha$-Potenz als Ordnungsprinzip ist nicht sichtbar. Die $\alpha=1$-Setzung
legt das frei.

\section{Warum diese Setzung -- Vergleich mit Alternativen}

\textbf{[S]} \textbf{Gegenüber Gauß ($e^2=\alpha$).} Funktioniert ebenfalls,
setzt die $4\pi$ aber in die Feldgleichungen statt ins Coulomb-Gesetz.
Heaviside-Lorentz setzt die $4\pi$ dorthin, wo sie geometrisch hingehören
(Kugeloberfläche, Coulomb), und lässt die Maxwell-Gleichungen $4\pi$-frei.
$\alpha$ ist dann die reine Kopplung, nicht ein Mischterm aus Kopplung und
Geometrie.

\textbf{[S]} \textbf{Gegenüber $c=\hbar=1$-Einheitenwahl.} Die Setzung
$c=\hbar=1$ vernichtet die Potenzen -- die Spur ist weg, die Rückrechnung
läuft nur noch über Dimensionsanalyse (A266). Die Definitionsgleichung
$e^2=4\pi\alpha$ dagegen \emph{behält die Spur}: die Kopplungspotenz ist
mitgeführt, weil die Relation stehen bleibt. Stipulation mit erhaltener Spur,
nicht Norm mit gelöschter -- deshalb konstruktiv reversibel, nicht bloß
begründet rekonstruierbar.

\section{FFGFT-Bezug: $\alpha$ zurückrechnen heißt $\xi$ einsetzen}

\textbf{[B]} $\alpha$ ist im FFGFT-Korpus keine freie Eingabe, sondern folgt
aus der Geometrie (A130):
\[
  \alpha^{-1} = 137{,}036 \quad\text{aus}\quad \xi = \frac{4}{30000},
  \quad D_f = 3-\xi, \quad E_0=\sqrt{m_em_\mu}.
\]
Die Setzung $\alpha=1$ isoliert genau die Stelle, an der die $\xi$-Geometrie
eingeht. Jedes zurückgesetzte $\alpha^n$ ist ein Ort, an dem $D_f=3-\xi$
sichtbar wird.

\textbf{[B]} Korpusintern heißt „$\alpha$ zurückrechnen'' deshalb nicht
„$1/137$ einsetzen'', sondern „die $\xi$-Struktur einsetzen'' -- und das
$\alpha{=}1$-System macht abzählbar, an wie vielen Stellen und in welcher
Potenz das geschieht.

\section{Zeichenerklärung}

Symbole die in der gesamten A-Serie verwendet werden, sind in A010 erklärt.
Spezifisch für dieses Dokument:

\begin{center}
\renewcommand{\arraystretch}{1.3}
{\small\begin{tabular}{lp{9cm}}
\toprule
Symbol & Bedeutung \\
\midrule
$e^2 = 4\pi\alpha$ & Heaviside-Lorentz-Definitionsgleichung \\
$e = \sqrt{4\pi}\approx3{,}545$ & Elementarladung bei $\alpha=1$ \\
$r_e = \alpha/m_e$ & Klassischer Elektronradius \\
$\lambda_C = 1/m_e$ & Compton-Wellenlänge \\
$a_0 = 1/(\alpha m_e)$ & Bohr-Radius \\
$V$ & Anzahl der Feynman-Vertizes im Prozess \\
$\alpha^n$ & $\alpha$-Ordnung / Rückfaktor des Prozesses \\
\bottomrule
\end{tabular}}
\renewcommand{\arraystretch}{1.0}
\end{center}

\section{Prüfskript}

\begin{itemize}[leftmargin=2em,itemsep=2pt,topsep=2pt]
  \item Numerische Verifikation aller Tabellenzeilen mit $\alpha=1/137{,}036$
    und $m_e=0{,}511\,\text{MeV}$; Verhältnis $r_e:\lambda_C:a_0=\alpha:1:\alpha^{-1}$:
    \texttt{python/A\_Serie\_Skripte/a267\_alpha\_eins.py}
\end{itemize}

\section{Was offen bleibt}

\textbf{Die Setzung $\alpha=1$ ist nur für elektromagnetische Prozesse
vollständig ausgearbeitet.} Für schwache und starke Kopplung ($\alpha_W$,
$\alpha_s$) wäre eine analoge Tabelle möglich, ist aber nicht ausgeführt.

\textbf{Der Zusammenhang $\alpha^n$ als Ort von $D_f=3-\xi$} ist für
einfache QED-Prozesse plausibel gezeigt, aber nicht für Mehrsschleifen-
oder nicht-abelsche Prozesse ausgeführt.

\section{Ersetzt}

Dieses Dokument nimmt auf: Dok.~011 (Einzelformeln $r_e$, $\lambda_C$, $a_0$,
Dimensionsanalyse $\alpha=r_e/\lambda_C$, SI-Zahlenwerte der atomaren Skalen),
Dok.~015 (Heaviside-Lorentz-Zeile $e^2/(4\pi)=1$ in der Einheitensystematik).
Die zusammenhängende Umrechnungstabelle mit expliziten $\alpha^n$-Rückfaktoren
und die 137-Leiter $r_e:\lambda_C:a_0=\alpha:1:\alpha^{-1}$ als Strukturaussage
sind im Altkorpus nicht als Einheit vorhanden und hier neu systematisiert.
Verwandt: A130 (geometrische Herleitung $\alpha$ aus $\xi$), A266 (SI-Rückrechnung).

\section{Verwendung von KI-Werkzeugen}

Dieses Dokument ist unter Verwendung KI-gestützter Werkzeuge entstanden.
Verantwortung für Inhalt und Fehler liegt beim Autor. Wortlaut der Regeln in A010.
